
\section*{\underline{Exposé VI}\\[3pt] \underline{\underline{LE @@M0@@ D'UN FIBRÉ PROJECTIF :}}\\[3pt] \underline{\underline{CALCULS ET CONSÉQUENCES}}}
% Exposé VI: the print's doubly underlined S and underlined Proj (scan check, 10 Nov)
\renewcommand{\Sym}{\mathbf{S}}
\renewcommand{\Projz}{\underline{\operatorname{Proj}}_0}
\begin{center}
\underline{par Pierre BERTHELOT}
\end{center}

\section*{1. Calcul du @@M1@@ d'un fibré projectif : cas des faisceaux localement libres de type fini}

Dans ce paragraphe, la notation @@M2@@ désignera le groupe de Grothendieck des faisceaux localement libres de type fini sur un schéma @@M3@@.

Le but du paragraphe est de démontrer le théorème suivant :

\textbf{Théorème 1.1.} \uline{Soient} @@M4@@ \uline{un schéma quasi-compact}, @@M5@@ \uline{un} @@M6@@\uline{-Module localement libre de rang} @@M7@@, @@M8@@ \uline{le fibré projectif associé} ; \uline{le morphisme structural} : @@M9@@ \uline{définit un homomorphisme d'anneaux} : @@M10@@ \uline{munissant} @@M11@@ \uline{d'une structure de} @@M12@@-\uline{algèbre}. \uline{Soit} @@M13@@ \uline{la classe du faisceau} @@M14@@. \uline{Alors} @@M15@@ \uline{est un} @@M16@@-\uline{module libre admettant pour base} @@M17@@.

\uline{Plan de la démonstration} :

@@M18@@) Les puissances de @@M19@@ engendrent @@M20@@ comme @@M21@@-module.

Pour le montrer, nous allons montrer les points suivants :

\begin{enumerate}[label=\alph*)]
\item Si @@M22@@ est un @@M23@@-Module localement libre de type fini, il existe un @@M24@@-Module gradué @@M25@@ de présentation finie, plat sur @@M26@@, tel que @@M27@@ (@@M28@@ étant l'algèbre symétrique de @@M29@@~; pour la notation @@M30@@, voir EGA II 8.12.1).

\item Il existe une résolution graduée finie de @@M31@@ par des @@M32@@-Modules gradués @@M33@@ localement libres de type fini : @@M34@@.

\item Les @@M35@@ admettent une filtration finie dont les quotients sont du type @@M36@@, les @@M37@@ étant localement libres de type fini sur @@M38@@.
\end{enumerate}

